

\section{Useful formulae}

\subsection{12d Gravity}
Written in Einstein frame, and real torus coordinates, with convention $(p,q)$ denoting (NSNS, RR) charges, the 12d metric is
\begin{equation}
\hat{g}_{MN}
= \begin{pmatrix}
    g_{mn} & 0 & 0 \\
    0 & \frac{1}{\tau_2} & \frac{\tau_1}{\tau_2} \\
    0 & \frac{\tau_1}{\tau_2} & \frac{|\tau|^2}{\tau_2}
\end{pmatrix},\quad
\hat{g}_{ab}
= \frac{1}{\tau_2}
\begin{pmatrix}
    1 & \tau_1 \\
    \tau_1 & |\tau|^2    
\end{pmatrix}
\end{equation}
From this 12d metric, one finds that the nonvanishing Christoffel symbols are $\hat{\Gamma}^p{}_{mn} = \Gamma^p{}_{mn}$, $\hat{\Gamma}^a{}_{bm} = \frac{1}{2}g^{ac}g_{bc,m}$, $\hat{\Gamma}^m{}_{ab} = - \frac{1}{2}g^{mp}g_{ab,p}$.
With $a$ indexing rows and $b$ columns:
\begin{equation}
\begin{aligned}
\Gamma^a{}_{bm}
&=\frac{1}{2 \tau_2^2}\begin{pmatrix}
    -\tfrac12 \nabla_m|\tau|^2 & |\tau|^2\nabla_m\tau_1 - \tau_1 \nabla_m|\tau|^2 \\
    \nabla_m\tau_1 & \tfrac12 \nabla_m|\tau|^2
\end{pmatrix}\\
\Gamma_{abm}=-\Gamma^m{}_{ab}
&= \frac{1}{2\tau_2}\begin{pmatrix}
    -\nabla_m\ln\tau_2
    & \nabla_m\tau_1 - \tau_1\nabla_m\ln\tau_2 \\
    \nabla_m\tau_1 - \tau_1\nabla_m\ln\tau_2 & \nabla_m|\tau|^2 - |\tau|^2\nabla_m\ln\tau_2
\end{pmatrix}
=\frac{1}{2}g_{ab,m}
\end{aligned}
\end{equation}
The components of $R_{manb}$ are given by

\medskip
\renewcommand{\arraystretch}{1.8}
\begin{tabular}{c|cccccc}
\toprule
$R_{ambn}$ 
  & $\nabla_m\nabla_n\tau_1$ 
  & $\nabla_m\nabla_n\tau_2$ 
  & $(\nabla_m\tau_1)(\nabla_n\tau_1)$ 
  & $\nabla_m\tau_1\nabla_n\tau_2$ 
  & $\nabla_m\tau_2\nabla_n\tau_1$ 
  & $(\nabla_m\tau_2)(\nabla_n\tau_2)$ \\
\midrule
$R_{1m1n}$ & $0$ & $\frac{1}{2\tau_2^2}$ & $\frac{1}{4\tau_2^3}$ & $0$ & $0$ & $-\frac{3}{4\tau_2^3}$ \\
$R_{1m2n}$ & $-\frac{1}{2\tau_2}$ & $\frac{\tau_1}{2\tau_2^2}$ & $\frac{\tau_1}{4\tau_2^3}$ & $\frac{1}{4\tau_2^2}$ & $\frac{3}{4\tau_2^2}$ & $-\frac{3\tau_1}{4\tau_2^3}$ \\
$R_{2m1n}$ & $-\frac{1}{2\tau_2}$ & $\frac{\tau_1}{2\tau_2^2}$ & $\frac{\tau_1}{4\tau_2^3}$ & $\frac{3}{4\tau_2^2}$ & $\frac{1}{4\tau_2^2}$ & $-\frac{3\tau_1}{4\tau_2^3}$ \\
$R_{2m2n}$ & $-\frac{\tau_1}{\tau_2}$ & $\frac{\tau_1^2-\tau_2^2}{2\tau_2^2}$ & $\frac{\tau_1^2-3\tau_2^2}{4\tau_2^3}$ & $\frac{\tau_1}{\tau_2^2}$ & $\frac{\tau_1}{\tau_2^2}$ & $\frac{\tau_2^2-3\tau_1^2}{4\tau_2^3}$ \\
\bottomrule
\end{tabular}
\medskip

% \renewcommand{\arraystretch}{1.8}
% \begin{tabular}{c|ccccc}
% \toprule
% $g^{mn}R_{ambn}$
%   & $\Delta\tau_1$
%   & $\Delta\tau_2$
%   & $(\nabla\tau_1)^2$
%   & $\nabla\tau_1\!\cdot\!\nabla\tau_2$
%   & $(\nabla\tau_2)^2$ \\
% \midrule
% $g^{mn}R_{1m1n}$ & $0$ & $\frac{1}{2\tau_2^2}$ & $\frac{1}{4\tau_2^3}$ & $0$ & $-\frac{3}{4\tau_2^3}$ \\[4pt]
% $g^{mn}R_{1m2n}$ & $-\frac{1}{2\tau_2}$ & $\frac{\tau_1}{2\tau_2^2}$ & $\frac{\tau_1}{4\tau_2^3}$ & $\frac{1}{\tau_2^2}$ & $-\frac{3\tau_1}{4\tau_2^3}$ \\[4pt]
% $g^{mn}R_{2m1n}$ & $-\frac{1}{2\tau_2}$ & $\frac{\tau_1}{2\tau_2^2}$ & $\frac{\tau_1}{4\tau_2^3}$ & $\frac{1}{\tau_2^2}$ & $-\frac{3\tau_1}{4\tau_2^3}$ \\[4pt]
% $g^{mn}R_{2m2n}$ & $-\frac{\tau_1}{\tau_2}$ & $\frac{\tau_1^2-\tau_2^2}{2\tau_2^2}$ & $\frac{\tau_1^2-3\tau_2^2}{4\tau_2^3}$ & $\frac{2\tau_1}{\tau_2^2}$ & $\frac{\tau_2^2-3\tau_1^2}{4\tau_2^3}$ \\
% \bottomrule
% \end{tabular}

We further find the trace is $g^{ab} R_{ambn} = -\frac{1}{2\tau_2^2}(\nabla_m \tau_1 \nabla_n \tau_1 + \nabla_m \tau_2 \nabla_n\tau_2)$.
The antisymmetric part is $\frac{1}{2}(R_{1m2n} - R_{2m1n})
= -\frac{1}{4\tau_2^2}(\nabla_m\tau_1 \nabla_n\tau_2 - \nabla_m\tau_2 \nabla_n\tau_1)
$.
We find
$$
R_{mn12} = \frac{1}{2\tau_2^2}[\nabla_m\tau_2\nabla_n\tau_1-\nabla_m\tau_1\nabla_n\tau_2]
$$
One can then go on to compute that
$$
R_{1212} = \frac{1}{4\tau_2^2}|\nabla \tau|^2.
$$
We note that generically $R_{ambn}$, $R_{1212}$ and $R_{mn12}$ are independent of each other, but here we see $R_{1212}$ and $R_{mn12}$ can both be obtained from linear combinations of $R_{ambn}$.
So the 12d Riemann tensors only provide 4 independent degrees of freedom for second order derivatives of $\tau$, while in total there are 6 independent degrees of freedom for second order derivatives of $\tau$, as given in the columns of the above table.


\subsection{IIB useful equations}

Reference: \cite{Policastro:2008hg,Minasian:2015bxa}

The rank-4 tensor $DP$ is defined by
\begin{equation}
(\nabla \mathcal{P})_{mn}{}^{pq} \equiv 
\delta_{[m}^{[p} \nabla_{n]}\mathcal{P}^{q]}
\end{equation}
which has the same symmetries as the Riemann tensor.
The $U(1)$ connection and one-form are defined by
\begin{equation}
Q_m =-\frac{\nabla_m\tau_1}{2\tau_2},\quad
P_m = \frac{i}{2\tau_2} \nabla_m \tau,
\end{equation}
From here one can compute
\begin{equation}
\begin{aligned}
D_m \mathcal{P}_n
&= \nabla_m \mathcal{P}_n - 2i Q_m \mathcal{P}_n
&= 2\mathcal{P}_m \mathcal{P}_n + \frac{i}{2\tau_2} \nabla_m \nabla_n \tau \\
D_m \bar{\mathcal{P}}_n
&= \nabla_m \bar{\mathcal{P}}_n + 2i Q_m \bar{\mathcal{P}}_n
&= 2\bar{\mathcal{P}}_m \bar{\mathcal{P}}_n - \frac{i}{2\tau_2} \nabla_m \nabla_n \bar{\tau}
\end{aligned}
\end{equation}
which implies that $D_m P_n$ is symmetric.
The complex 3-form is defined by
\begin{equation}
G_3 \equiv \frac{1}{\sqrt{\tau_2}} (F_3 - \tau H_3)
\end{equation}
The contraction of the $t_8 t_8$ tensor with four rank-4 tensors is defined by
\begin{align}
	t_8 t_8 \left\{ ABCD \right\} := (t_8)^{a_1 \ldots a_8} (t_8)^{b_1 \ldots b_8} A_{a_1 a_2 b_1 b_2} B_{a_3 a_4 b_3 b_4} C_{a_5 a_6 b_5 b_6} D_{a_7 a_8 b_7 b_8} \,,
\end{align}
The $t_8$ tensor itself is conventionally \cite{Aggarwal:2025lxf,Liu:2025uqu} defined to represent taking traces with cyclic symmetry:
\begin{equation}
t_8 ABCD = 8 \text{tr} (ABCD) - 2 \text{tr} (AB) \text{tr} (CD) + \text{cycl}(B,C,D).
\end{equation}
It also has a component definition \footnote{See also \cite{Minasian:2015bxa}}
\begin{equation}
\begin{aligned}
t_8^{a_1 a_2 b_1 b_2 c_1 c_2 d_1 d_2}
&= 8 g^{a_2 b_1} g^{b_2 c_1} g^{c_2 d_1} g^{d_2 a_1}
- 2 g^{a_2 b_1} g^{b_2 a_1} g^{c_2 d_1} g^{d_2 c_1}\\
&+ \text{cycl}(b,c,d)
- (a_1 \leftrightarrow a_2)
- (b_1 \leftrightarrow b_2)
- (c_1 \leftrightarrow c_2)
- (d_1 \leftrightarrow d_2).
\end{aligned}
\end{equation}
In particular, when acting on a tensor with the symmetries of the Riemann tensor, we have

\begin{align}
t_8 t_8 R^4 &= 3 \times 2^7 \Biggl\{
R_{abef} R^{fg}{}_{cd} R^{dahe} R_{gh}{}^{bc}
+ \frac{1}{2} R_{abef} R^{fgbc} R_{cdgh} R^{heda} \notag \\
&\quad - \frac{1}{2} R_{abef} R^{efbc} R_{cdgh} R^{dagh}
- \frac{1}{4} R_{abef} R^{ef}{}_{cd} R^{dagh} R_{gh}{}^{bc} \notag \\
&\quad + \frac{1}{16} R_{abef} R^{abgh} R_{ghcd} R^{cdef}
+ \frac{1}{32} (R_{abcd} R^{abcd})^2
\Biggr\}. \label{eq:A2}
\end{align}
Another useful identity is if $S$ has Riemann-like symmetries, then
\begin{align}
t_8 t_8 R^2 S^2 &= 4 R_{abcd} R_{abcd} S_{efgh} S_{efgh}
+ 16 R_{abcd} R_{abef} S_{cdgh} S_{efgh} \\
&\quad + 8 R_{abcd} R_{efgh} S_{abcd} S_{efgh}
+ 8 R_{abcd} R_{efgh} S_{abef} S_{cdgh} \\
&\quad - 64 R_{abcd} R_{abce} S_{dfgh} S_{efgh}
- 32 R_{abcd} R_{abef} S_{cegh} S_{dfgh} \\
&\quad - 64 R_{abcd} R_{aefg} S_{bhfg} S_{cdeh}
- 64 R_{abcd} R_{efgh} S_{abce} S_{dfgh} \\
&\quad - 64 R_{abcd} R_{aefg} S_{bhcd} S_{ehfg}
+ 128 R_{abcd} R_{aecf} S_{bgdh} S_{egfh} \\
&\quad + 256 R_{abcd} R_{aefg} S_{bhcf} S_{dgeh}
+ 128 R_{abcd} R_{aecf} S_{bgfh} S_{dheg} \\
&\quad + 64 R_{abcd} R_{efgh} S_{aecg} S_{bfdh}, 
\label{eq:t8R2S2}
\end{align}
